% s-023-p023-m002; original PDF 23, 24
\Aphorism{s-023-p023-m002}{14}
Von einer bestimmten Complicationsstufe des menschlichen Trieblebens an ist die Beziehung zwischen Wille und Handlung, Empfindung und Ausdruck, Charakter und Bethätigung nicht mehr die der Übereinstimmung, sondern die des Widerspruchs; ein Affect, der über die ersten unschuldigen Geh- und Sprechversuche hinaus ist und zu raffiniren anfängt, scheint in der Darstellung nach aussen das antimechanische Princip zu befolgen: möglichst viel Unlust, möglichst starke Widerstände, möglichst grosse Kraftverschwendung! So äussert sich dann Wohlwollen als Härte und Grausamkeit, Empfindlichkeit für Leid und Mitleid als frivole Leugnung des Leidens, Tiefe des Erkennens als Oberflächlichkeit, Machtbegier als Unterwerfung, Gefühlsüberschwang als plötzliche Kälte und Nüchternheit, Sinnlichkeit als Askese, Rachsucht als Grossmuth, Glaubensleidenschaft als gewaltsamer Materialismus — — und wie jene seltenen und prachtvollen Affect-Entladungen ins Gegensätzliche alle heissen mögen, bei denen ein ganzer Wolken- und Gewitterhimmel mit Einem Schlage zu taghellem Feuerschein aufflammt. Mit solchen Sturmfluthen, Deichbrüchen, Kraftexplosionen weiss die normale Psychologie nichts anzufangen; diese \Sperr{Heteropathie} der Seele tritt auf, wenn die ältere, regelrechte Form von Affectäusserung, die \Sperr{homopathische} (wo Gleiches immer Gleiches bedeutet) einmal für eine Zeit lang zu den abgelaufenen Heerstrassen und ausgefahrenen Geleisen gehört. Als dritte Form, die aber schon einer vollkommenen Beherrschung und Verstellung des Trieblebens gleichkommt, wäre eine Relation zwischen Innen und Aussen denkbar, die weder einfach abschreibt noch einfach widerspricht, sondern bald das Eine, bald das Andere thut und in der zufälligen Vertheilung der Fälle jeden Rückschluss von Aussen nach Innen entkräftet.
\par

